% ============================================================
% Introduction and Related Work — Draft
% Paper: Do Morphological Auxiliary Losses Actually Improve Crack Segmentation?
%        A Systematic Multi-Seed Study
% Target: Automation in Construction (elsarticle-num)
% ============================================================

\section{Introduction}
\label{sec:introduction}

Automated detection and quantification of surface defects---primarily cracks and spalling---is a critical task in structural health monitoring of concrete infrastructure~\cite{koch2015review, mosalam2024dlshm}. Manual visual inspection remains labour-intensive, subjective, and difficult to scale across geographically distributed assets such as dams, bridges, and tunnels. Over the past decade, deep learning has transformed this domain: convolutional and transformer-based semantic segmentation models now routinely achieve high per-pixel accuracy on benchmark crack datasets~\cite{kulkarni2022crackseg9k, dorafshan2018sdnet, zhang2025crackreview}.

The dominant paradigm employs encoder--decoder architectures---U-Net~\cite{ronneberger2015unet}, DeepLabV3+~\cite{chen2018deeplabv3plus}, and more recently SegFormer~\cite{xie2021segformer}---trained with standard pixel-level losses such as cross-entropy and Dice. While these models perform well on in-distribution data, two persistent challenges motivate ongoing research. First, cracks are thin, elongated structures whose \emph{topological} integrity (connectivity, branching) is poorly captured by area-based loss functions~\cite{shit2021cldice, hu2019topology}. This has led to a growing body of work on topology-preserving auxiliary losses and morphological supervision heads that aim to improve segmentation quality beyond what per-pixel objectives can achieve. Second, models trained on imagery from one inspection site often generalise poorly to other sites, sensors, or surface conditions---a cross-site domain gap that remains largely unaddressed in the structural inspection literature~\cite{bukhsh2021damage, hassani2025dashm}.

To tackle the first challenge, researchers have proposed auxiliary losses and supervision signals drawn from crack morphology: centerline Dice (clDice)~\cite{shit2021cldice}, persistent-homology-based topological losses~\cite{clough2020topological, hu2019topology}, skeleton distance transform regression~\cite{audebert2019distance}, and keypoint heatmaps for endpoints and junctions~\cite{sun2019crackjunction}. These methods are widely assumed to improve segmentation quality, and published results often show gains of one to several percentage points in mean Intersection-over-Union (mIoU) or similar metrics. However, a critical methodological concern undermines confidence in these findings: the vast majority of studies report results from a \emph{single random seed}, making it impossible to distinguish genuine methodological improvements from seed-dependent variance~\cite{picard2021seed, henderson2018rl}. Recent work has demonstrated that random seed selection alone can produce performance differences comparable in magnitude to claimed improvements, both in computer vision generally~\cite{picard2021seed} and in civil engineering applications specifically~\cite{alsabbagh2024randomness}.

For the second challenge, standard domain generalisation (DG) methods---including CORAL~\cite{sun2016coral}, DANN~\cite{ganin2016dann}, and MixStyle~\cite{zhou2021mixstyle}---have been developed to learn domain-invariant representations. While these methods have demonstrated effectiveness in multi-source settings (e.g., photo/art/sketch domains), their applicability to single-source infrastructure inspection, where pseudo-domains must be constructed from within-dataset partitions, has not been systematically evaluated.

This paper presents a rigorous, multi-seed empirical study that challenges both lines of assumption. We construct a progressive baseline ladder on SegFormer-B2~\cite{xie2021segformer} with a shared Feature Pyramid Network~\cite{lin2017fpn}, systematically adding morphological supervision---clDice, skeleton distance transform regression, keypoint heatmaps, and width estimation---and evaluating each configuration with pre-registered go/no-go gates. We further evaluate three standard DG methods under a pseudo-domain setting using difficulty-based data partitioning. All experiments are validated with three independent random seeds, and statistical significance is assessed via paired $t$-tests.

Our central finding is \textbf{negative}: no form of auxiliary morphological supervision significantly improves crack segmentation over the mask-only baseline when validated with multiple seeds, and no standard DG method reduces the cross-site domain gap under pseudo-domain single-source training. More critically, we identify \textbf{three cases of misleading single-seed results} that were overturned by multi-seed validation:
\begin{enumerate}
    \item Skeleton distance transform supervision appeared to improve mIoU\textsubscript{fg} by +1.0\% in a single seed, but multi-seed evaluation revealed a non-significant +0.5\%$\pm$0.3\% ($p=0.145$).
    \item clDice appeared to degrade performance by $-$1.6\% in a single seed, but multi-seed evaluation showed a non-significant +0.2\%$\pm$0.9\% ($p=0.822$).
    \item MixStyle appeared to improve out-of-distribution performance by +52\% in a single seed, but multi-seed evaluation showed statistical equivalence to the ERM baseline ($p=0.878$).
\end{enumerate}

These results carry implications beyond our specific experimental setting. They suggest that the crack segmentation literature may contain a non-trivial number of fragile claims built on single-seed comparisons, and they underscore the need for multi-seed validation as a minimum standard of evidence for claimed improvements.

The contributions of this work are as follows:
\begin{itemize}
    \item A systematic evaluation of five auxiliary morphological supervision strategies (clDice, skeleton distance transform, endpoint/junction detection, width regression, and their combination) for concrete defect segmentation, demonstrating that none significantly improves over a mask-only baseline under multi-seed validation.
    \item A multi-seed evaluation of three standard domain generalisation methods (CORAL, DANN, MixStyle) under pseudo-domain single-source training, showing that all are statistically equivalent to empirical risk minimisation.
    \item Identification of three concrete cases where single-seed evaluation produced misleading conclusions, providing an empirical demonstration of the fragility of single-seed comparisons in this domain.
    \item A multi-tier graph evaluation protocol for assessing the quality of crack graphs extracted from segmentation predictions, with strict, relaxed, and lenient tolerance tiers.
\end{itemize}


\section{Related Work}
\label{sec:related_work}

We review three research areas that our systematic study directly evaluates: topology-preserving losses for crack segmentation, crack graph and skeleton extraction methods, and domain generalisation for infrastructure defect detection.

\subsection{Topology-Preserving Losses for Crack Segmentation}
\label{sec:rw_topology}

Standard segmentation losses such as cross-entropy and Dice operate on a per-pixel basis and are agnostic to the topological structure of predicted regions. For thin, elongated objects like cracks, vessels, and roads, this can result in segmentations that achieve high volumetric overlap but contain topologically incorrect disconnections or spurious connections. This limitation has motivated a substantial body of work on topology-aware loss functions.

Hu et al.~\cite{hu2019topology} introduced a topological loss based on persistent homology that penalises differences in Betti numbers between predicted and ground-truth segmentations, with demonstrated improvements on neuron membrane and vessel segmentation. Clough et al.~\cite{clough2020topological} extended this approach with a differentiable persistent homology loss applicable to both 2D and 3D segmentation tasks. While theoretically principled, persistent-homology-based losses are computationally expensive and can be difficult to combine with standard losses during training~\cite{hu2021topology}.

An alternative line of work exploits morphological skeletons rather than algebraic topology. Shit et al.~\cite{shit2021cldice} proposed the centerline Dice (clDice) loss, which computes Dice overlap on soft morphological skeletons of the prediction and ground truth. The clDice loss is computationally efficient and theoretically guarantees topology preservation up to homotopy equivalence for binary segmentation. It has since been widely adopted for vessel, road, and crack segmentation. Kirchhoff et al.~\cite{kirchhoff2024skeleton} proposed the Skeleton Recall Loss, which performs tubed skeletonisation on the ground truth and computes a soft recall loss, reducing computational overhead by over 90\% compared to clDice while achieving competitive or superior connectivity preservation.

In the crack detection domain specifically, He and Li~\cite{he2024beyond} evaluated skeleton-based losses for preserving connectivity in crack segmentation, demonstrating improvements on the EdmCrack600 and CrackForest datasets using U-Net++ and SegFormer backbones. Jing et al.~\cite{jing2025topology} proposed TopoM-CrackNet, which combines persistent homology with a Mamba-enhanced U-Net for pavement crack detection with topology preservation, published in \textit{Automation in Construction}. Zhang and Qi~\cite{zhang2024irfusionformer} introduced a topology-based loss within an RGB-thermal fusion framework for pavement crack segmentation. Xu et al.~\cite{xu2021cploss} proposed CP-loss for connectivity preservation in road curb detection, which focuses on regions of disconnectivity during optimisation.

A critical observation across this literature is that the vast majority of studies report results from a single random seed. While the methods are often evaluated on multiple datasets, the lack of multi-seed validation within each dataset makes it difficult to assess whether reported improvements exceed the variance attributable to random initialisation. Our study directly addresses this gap by evaluating clDice~\cite{shit2021cldice} under multi-seed conditions and finding that its effect on crack segmentation is statistically indistinguishable from zero.

\subsection{Crack Graph and Skeleton Extraction}
\label{sec:rw_graph}

Beyond pixel-level segmentation, practical crack assessment requires measurable morphological features: centerline geometry, branching topology, junction locations, crack width, and crack length~\cite{li2024crackquantification, lin2025crackmorph}. Extracting a graph representation---where nodes represent endpoints and junctions, and edges represent crack segments---from a segmentation mask is a natural approach to obtaining such features.

Traditional approaches rely on morphological skeletonisation of binary crack masks, followed by graph extraction from the resulting skeleton pixels~\cite{wang2025crackskeletonization}. Endpoints (degree-1 pixels) and junctions (degree$\geq$3 pixels) are identified directly from the skeleton topology, and connected skeleton segments between these keypoints form graph edges. This pipeline is straightforward but sensitive to segmentation noise: small errors in the predicted mask can produce dramatically different skeletons, leading to fragile graph representations. Li and Yuan~\cite{li2025cracktopology} recently demonstrated a graph-based quantification framework for automated topological analysis of crack networks in road maintenance.

Distance transform (DT) maps provide an alternative representation that encodes the distance of each pixel to the nearest boundary~\cite{audebert2019distance}. Ridge detection on DT maps can recover crack centerlines without explicit skeletonisation, and local maxima correspond to the widest points of crack segments. Audebert et al.~\cite{audebert2019distance} showed that auxiliary DT regression can spatially regularise semantic segmentation predictions, and this approach has been adopted in various infrastructure inspection contexts.

Junction detection in crack images has been addressed through correlation structure analysis and tensor voting~\cite{sun2019crackjunction}, achieving high correctness and completeness on pavement images. More recently, Liu and Li~\cite{liu2024ltggnet} proposed LTGGNet, which uses local topology and global grouping attention for concrete crack detection. Lin and Zheng~\cite{lin2025crackmorph} introduced CrackMorph-XAI-Net, an explainable framework that integrates topology-preserving skeleton extraction with junction detection via Gaussian heatmap regression, producing measurable morphological descriptors. Kuttner and Grunwald~\cite{kuttner2025cracktopo} proposed extensions to the crack topology score for evaluating the topological quality of crack segmentations.

A common assumption in this literature is that auxiliary supervision targeting morphological features---skeleton DT regression, keypoint detection---should improve the quality of downstream graph extraction. Our study tests this assumption empirically: we find that skeleton DT supervision (B2) does not significantly improve post-hoc graph extraction quality compared to a mask-only baseline (B0), with all Wilcoxon signed-rank tests yielding $p > 0.23$. We further find that direct neural graph prediction using CenterNet-style~\cite{zhou2019centernet} node heatmaps and edge classifiers, while architecturally sound (demonstrated by overfit gate passage), is limited by the quality of automatically generated training labels.

\subsection{Domain Generalisation for Infrastructure Defect Detection}
\label{sec:rw_dg}

Deep learning models for structural inspection are typically trained on data from a single site or a small number of sites and are expected to generalise to unseen infrastructure. In practice, substantial performance degradation is observed when models are deployed on imagery from different sites, sensors, or environmental conditions~\cite{koch2015review, bukhsh2021damage}. This cross-site domain gap is well documented: Zhang and Wei~\cite{zhang2025crackreview} note that generalisability across datasets remains one of the most significant open challenges in crack detection. Hassani and Dackermann~\cite{hassani2025dashm} provide a systematic review of domain adaptation methods for structural health monitoring, covering both vibration-based and vision-based approaches.

Domain adaptation (DA) methods, which assume access to unlabelled target data during training, have been applied to crack segmentation with some success. Jiang and Yuan~\cite{jiang2025udacrack} proposed a cross-domain stylisation framework with dual adversarial feature learning for crack segmentation across different surface types. Jiang and Deng~\cite{jiang2025crossdataset} demonstrated cross-dataset semantic segmentation for composite crack detection using unsupervised transfer learning. Bukhsh et al.~\cite{bukhsh2021damage} evaluated in-domain and cross-domain transfer learning for structural damage detection, finding that cross-domain transfer is substantially more challenging.

Domain generalisation (DG) methods, which aim to learn domain-invariant representations \emph{without} access to target data, represent a more challenging but practically relevant setting. Three widely used DG methods are: (1)~Deep CORAL~\cite{sun2016coral}, which aligns second-order feature statistics across domains; (2)~DANN~\cite{ganin2016dann}, which uses a gradient reversal layer to learn features that are discriminative for the task but invariant to domain identity; and (3)~MixStyle~\cite{zhou2021mixstyle, zhou2023mixstyle}, which synthesises new domains by mixing instance-level feature statistics during training. These methods have been developed and validated primarily in multi-source DG benchmarks (e.g., PACS, Office-Home) where distinct visual domains are available.

A key practical constraint in infrastructure inspection is that labelled data often comes from a single source domain. In this single-source DG setting, pseudo-domains must be constructed---for example, by partitioning data based on difficulty, augmentation style, or other heuristics~\cite{elghazy2025pseudodg}. However, the effectiveness of pseudo-domain construction is not guaranteed: if the constructed pseudo-domains do not capture genuine distributional variation, DG methods may fail to learn meaningful domain-invariant features. Elghazy and Alaa~\cite{elghazy2025pseudodg} note that if a pseudo-domain's distribution is semantically irrelevant, it can disrupt the learning process and degrade model performance.

Our study evaluates all three DG methods (CORAL, DANN, MixStyle) in a single-source pseudo-domain setting where DamSegment training images are partitioned into difficulty-based tiers (Easy/Medium/Hard). Multi-seed evaluation with three seeds reveals that all three methods are statistically equivalent to the ERM baseline ($p=0.878$ for MixStyle vs.\ ERM on out-of-distribution mIoU\textsubscript{fg}), with a stable domain gap of approximately 55\%. This finding suggests that difficulty-based pseudo-domains do not capture the distributional characteristics needed for effective domain generalisation in concrete defect detection, and that bridging real cross-site gaps likely requires genuine multi-source data collection rather than algorithmic solutions applied to single-source data.
